%% ========================================================================
%% Appendix B: Solusi Latihan
%% ========================================================================

\chapter{Solusi Latihan}
\label{app:exercise-solutions}

Lampiran ini berisi solusi dan pembahasan untuk latihan-latihan di setiap bab.

%% ------------------------------------------------------------------------
\section{Solusi Bab 1: Pengenalan Sistem Operasi}

\subsection{Latihan 1.1}
\textbf{Soal:} Jelaskan perbedaan antara sistem operasi proprietary dan open-source.

\textbf{Jawaban:}
\begin{itemize}
    \item \textbf{Proprietary OS:} Sistem operasi yang source code-nya tertutup dan
          dimiliki oleh perusahaan tertentu. Contoh: Windows, macOS. Pengguna harus
          membayar lisensi dan tidak dapat memodifikasi source code.

    \item \textbf{Open-source OS:} Sistem operasi dengan source code yang terbuka
          untuk umum. Contoh: Linux, FreeBSD. Pengguna dapat melihat, memodifikasi,
          dan mendistribusikan ulang source code sesuai lisensi yang berlaku.
\end{itemize}

%% [Tambahkan solusi latihan lainnya]

%% ------------------------------------------------------------------------
\section{Solusi Bab 2: Manajemen Perangkat Keras}

%% [Konten akan diisi di sini]

%% ------------------------------------------------------------------------
\section{Solusi Bab 3: Dasar I/O}

%% [Konten akan diisi di sini]

%% [Lanjutkan untuk bab-bab lainnya...]
